% output=pdftex interface=en 

\environment man-init \mainlanguage[en]

\starttext

\titelblad
  {PPCH\TeX}
  {a macropackage for typesetting\\
   chemical structure formulas}
  {J. Hagen \& A.F. Otten\\
   Pragma ADE, Hasselt NL}

\title{Contents}

\setuppagenumber[state=stop] \placecontent \pagina

\title{Introduction}

\PPCHTEX\ is a set of coherent macros that can be used to typeset chemical
structure formulas. The macros fall back on \PICTEX, a public domain drawing
package written by Michael Wichura. Although originally written on top of
\PICTEX, the second release can cooperate with \PSTRICKS\ of Timothy Van
Zandt, that is, with some limitations. Personally we still prefer the quality
of the \PICTEX\ output. 

The macros can be used from within several macro packages and fall back on
a few generic \CONTEXT\ modules. The macros are written in a way that permits
relatively easy upward compatible extensions. The interface is conform the
\CONTEXT\ interface. 

\PPCHTEX\ was originally meant for typesetting chemical structure formulas
like sixrings. At the moment there is also support for reaction mechanisms.
Formulas can be typeset at different sizes. Common elements or frequently
used formula components can be reused. 

Flexibility, simplicity and quality have been prefered over speed. We don't
use the \PICTEX\ option to save pictures, because timing showed that the gain
was minimal. 

The first version of \PPCHTEX\ was ready for use in 1995. This manual
describes the second major release. Thanks to the many suggestions of Tobias
Burnus, Dirk Kuypers and Ton Otten, the functionality was extended 
considerably. 

\page \setuppagenumber[state=start]

\vervolgblad
  {Explanation}
  []
  [SIX,B,C,R1356,RZ1356]
  [NO_2,NO_2,NO_2,CH_3]

\chapter{Structures}

The number of commands that is used to typeset chemical structure formulas
is, apart from some bell and whistle commands, limited to four.\footnote{The
concept structure in this manual only refers to the chemical structure. It is
not related to the document structure.} In the following example all of these
commands are used. 

\startbuffer
\setupchemical[axis=on,frame=on]

\startchemical
  \chemical[SIX,B,R,RZ][1,2,3,4,5,6]
\stopchemical
\stopbuffer

\startfiguretext
  [left][]
  {none}
  {\getbuffer}
\voorbeeld[ex:sixring]
\typebuffer
\stopfiguretext

With \type{\setupchemical} we can influence the makeup of the formulas. These
setups influence all the following formulas, unless they are superceded by
local setup variables.\footnote{One can of course limit the scope of the
variables by using \type{{ }} and/or the grouping macros \type{..group}.} 

The set up variables can be defined right after \tex{startchemical}. In that
case the set up is only applied to one structure formula. 

\startbuffer
\startchemical[frame=on,width=fit,height=fit]
  \chemical[CARBON,CB1][A,B,C,D,E,F]
\stopchemical
\stopbuffer

\startfiguretext
  [left][]
  {none}
  {\getbuffer}
\voorbeeld
\typebuffer
\stopfiguretext

Both examples show that \type{\chemical} is the essential command. This
command that may be used more than once within a
\type{\start}||\type{\stop}||pair, is accompanied with two arguments. These
arguments are written between~\type{[ ]}. The first argument is used for
defining the chemical bonds, the second argument for the atoms and molecules
that make up the structure. 

Text is typeset in mathematical mode, this means that you may type anything
that normally is allowed between~\type{$ $}. 

We will explain the first example in more detail. The key \type{SIX} means
that we want to draw a sixring. In analogy we could type \type{ONE},
\type{THREE}, \type{FOUR} and \type{FIVE}, \type{EIGHT}, \type{CARBON},
\type{NEWMAN}, \type{CHAIR}, some alternatives on these keys and some
symbols. 

\def\littlestructure[#1]%
  {\vbox
     {\startchemical[frame=on,scale=small,size=small]%
      \chemical[#1][1,2,3,4,5,6,7,8]%
      \stopchemical}}

\placefigure{none}
 \startcombination[4*3]
   {\littlestructure[ONE,SB,Z]}           {One}
   {\littlestructure[THREE,B,R,RZ]}       {Three}
   {\littlestructure[FOUR,B,R,RZ]}        {Four}
   {\littlestructure[FIVE,B,R,CRZ]}       {Five}
   {\littlestructure[SIX,B,R,RZ]}         {Six}
   {\littlestructure[EIGHT,B]}            {Eight}
   {\littlestructure[FIVE,FRONT,B,-R,+R]} {Five Front}
   {\littlestructure[SIX,FRONT,B,-R,+R]}  {Six Front}
   {\littlestructure[CARBON,CB]}          {Carbon}
   {\littlestructure[NEWMAN,STAGGER,CB]}  {Newman Stacker}
   {\littlestructure[NEWMAN,ECLIPSE,CB]}  {Newman Eclipse}
   {\littlestructure[CHAIR,B,-R,+R]}      {Chair}
 \stopcombination

The dimensions of \type{CHAIR} are somewhat different from the others. This
structure is also different in other means. Rotation for example is not
possible. 

\startbuffer
\startchemical[frame=on,scale=small]
  \chemical
    [CHAIR,B,+R,-R,
     +RZ1,+RZ2,+RZ3,+RZ4,+RZ5,+RZ6,
     -RZ1,-RZ2,-RZ3,-RZ4,-RZ5,-RZ6]
    [a,b,c,d,e,f,1,2,3,4,5,6]
\stopchemical
\stopbuffer

\startfiguretext
  [left][]
  {none}
  {\getbuffer}
\voorbeeld
\typebuffer
\stopfiguretext

Within a structure the chemical bonds between the \kap{c}||atoms are defined
in the same way. In this example we use~\type{B} and~\type{R}. Bonds within a
structure are numbered and can be defined by: 

\starttyping
\chemical[SIX,B1,B2,B3,B4,B5,B6]
\chemical[SIX,B135]
\chemical[SIX,B1..5]
\stoptyping

These keys draw parts of a sixring. With \type{R} and \type{RZ} we place
substituents on the ring. The key \type{R} draws the bond from a ring corner
to the substituent ($\angle~120^\circ$). The corner is also identified with a
number. 

\starttyping
\chemical[SIX,B1..6,R1..6]
\stoptyping

The defininition above draws the six bonds in the sixring and the bonds to
the substituents. The substituents are identified by the key \type{RZ}. Again
numbers are used to mark the position. The substituents themselves are
defined as text in the second argument. 

\startbuffer
\startchemical[frame=on,width=6000]
  \chemical[SIX,B1..6,R1..6,RZ1..3][CH_3,CH_3,OH]
\stopchemical
\stopbuffer

\startfiguretext
  [left][]
  {none}
  {\getbuffer}
\voorbeeld
\typebuffer
\stopfiguretext

When the second argument is left out no text (substituents) are placed on
the ring and the key \type{RZ1..3} has no effect. 

\chapter{Bonds}

This chapter gives an overview of the bonds you can use in structures. From
the examples throughout this manual the use of the different keys will become
more meaningful. 

Bonds always have two alternatives: a long and a short version. The shortened
bonds leave room to place atoms within a structure. A number of bonds can be
shortened on both sides left (\type{-}) or right (\type{+}). 

\placetable
{Single bonds.}
\starttable[|rT|lw(4cm)|rT|lw(4cm)|]
\HL
\VL ~~~B \VL Bond          \VL ~~SB \VL Single Bond       \VL\FR
\VL ~~BB \VL Bold Bond     \VL ~-SB \VL Left Single Bond  \VL\MR
\VL ~~HB \VL Hydrogen Bond \VL ~+SB \VL Right Single Bond \VL\LR
\HL
\stoptable

The example below shows a number of bonds combined within one structure: 

\startbuffer
\startchemical[frame=on]
  \chemical[ONE,SD1,SB4,BB2,SB7,Z01247][C,H,H,H,H]
\stopchemical
\stopbuffer

\startfiguretext
  [left][]
  {none}
  {\getbuffer}
\voorbeeld
\typebuffer
\stopfiguretext

A bond can be followed by one or more numbers or a range, for example:
\type{B1}, \type{B135} and \type{B1..5}. When you want to draw all bonds you
can type~\type{B}. 

Within a ring structure you can define extra bonds between atoms, for example
a double or triple bond. 

\placetable
{Multiple bonds.}
\starttable[|rT|lw(4cm)|rT|lw(4cm)|]
\HL
\VL ~~EB \VL Extra Bond    \VL ~~DB \VL Double Bond       \VL\FR
\VL      \VL               \VL ~~TB \VL Triple Bond       \VL\LR
\HL
\stoptable

Free electrons and electron pairs can be defined in different ways.
The accompanying keywords start with an~\type{E}. 

\placetable
{Free electrons and electron pairs.}
\starttable[|rT|lw(4cm)|rT|lw(4cm)|]
\HL
\VL ~~ES \VL Extra Single  \VL ~~ED \VL Extra Double \VL\FR
\VL ~~EP \VL Extra Pair    \VL ~~ET \VL Extra Triple \VL\LR
\HL
\stoptable

The example below shows a carbon atom with 8~outer electrons arranged in a
chemically very peculiar way. 

\startbuffer
\startchemical[frame=on]
  \chemical[ONE,Z0,ES1,ED3,ET5,EP7][C]
\stopchemical
\stopbuffer

\startfiguretext
  [left][]
  {none}
  {\getbuffer}
\voorbeeld
\typebuffer
\stopfiguretext

Within a ring structure you can make a shortcut from one atom to another. In
that case the atom that you want to skip has to be identified. As a
replacement of the double bonds in an aromatic sixring a circle can be drawn
and charges can be placed within the ring. 

\placetable
{Special bonds.}
\starttable[|rT|lw(4cm)|rT|lw(4cm)|]
\HL
\VL ~~SS \VL Short Shortcut       \VL ~~~S \VL Shortcut            \VL\FR
\VL ~-SS \VL Left Short Shortcut  \VL ~MID \VL Open Mid Shortcut   \VL\MR
\VL ~+SS \VL Right Short Shortcut \VL MIDS \VL Closed Mid Shortcut \VL\LR
\HL
\stoptable

\placetable
{Circle bonds.}
\starttable[|rT|lw(4cm)|rT|lw(4cm)|]
\HL
\VL ~~~C \VL Circle               \VL ~~CD \VL Dashed Circle          \VL\FR
\VL ~~CC \VL Shifted Circle       \VL ~CCD \VL Dashed Shifted Circle  \VL\LR
\HL
\stoptable

An example will explain the use of the circular bond and the use of displaced
charges. 

\startbuffer
\startchemical[frame=on]
  \chemical
    [SIX,B,ER6,CCD12346,Z0,PB:RZ6,ONE,SB8,EP6,Z0,ZT6,Z8,PE]
    [\ominus,O,\oplus,H]
\stopchemical
\stopbuffer

\startfiguretext
  [left][]
  {none}
  {\getbuffer}
\voorbeeld
\typebuffer
\stopfiguretext

Substituents can be connected to all corners of a structure. A substituent
can be anything you want. It depends on the presence of atoms or molecules
whether the bonds are long or short. In the examples you will see a great
number of the keys that are used to define substituents. 

\placetable
{Bonds to substituents.}
\starttable[|rT|lw(4cm)|rT|lw(4cm)|]
\HL
\VL ~~~R \VL Radical       \VL ~~SR \VL Single Radical       \VL\FR
\VL ~~-R \VL Left Radical  \VL ~-SR \VL Single Left Radical  \VL\MR
\VL ~~+R \VL Right Radical \VL ~+SR \VL Single Right Radical \VL\LR
\HL
\stoptable

There are a few alternatives to draw bridges. 

\placetable
{Special bonds to substituents.}
\starttable[|rT|lw(4cm)|rT|lw(4cm)|]
\HL
\VL ~~RD \VL Radical Dashed       \VL ~~RB\VL Radical Bold       \VL\FR
\VL ~-RD \VL Left Radical Dashed  \VL ~-RB\VL Left Radical Bold  \VL\MR
\VL ~+RD \VL Right Radical Dashed \VL ~+RB\VL Right Radical Bold \VL\LR
\HL
\stoptable

Radicals can be drawn in three ways.\footnote{The word radical schould not be
interpreted chemically, but typographically.} Some alternatives are seldom
used. 

\startbuffer
\startchemical[frame=on]
  \chemical[SIX,B,R14,RD25,RB36]
\stopchemical
\stopbuffer

\startfiguretext
  [left][]
  {none}
  {\getbuffer}
\voorbeeld
\typebuffer
\stopfiguretext

\placetable
{More special bonds to substituents.}
\starttable[|rT|lw(4cm)|rT|lw(4cm)|]
\HL
\VL ~~SD \VL Single Dashed \VL ~LDD \VL Left Double Dashed  \VL\FR
\VL ~~OE \VL Open Ended    \VL ~RDD \VL Right Double Dashed \VL\LR
\HL
\stoptable

An example of an {\em Open Ended} is defined below. We see a sixring
(\type{SIX}) with a number of consecutive \type{ONE}s. The use of \type{PB}
is explained later. 

\startbuffer
\startchemical
    [width=5000,top=2500,bottom=1500,frame=on]
  \chemical
    [SIX,B,C,R6,
     PB:RZ6,ONE,CZ0,OE1,SB5,MOV5,CZ0,OFF5,OE5,PE]
    [CH,CH_2]
\stopchemical
\stopbuffer

\startfiguretext
  [left][]
  {none}
  {\getbuffer}
\voorbeeld
\typebuffer
\stopfiguretext

It's obvious that substituents can be attached to the structure by means of 
double bonds. 

\placetable
{Double bonds to substituents.}
\starttable[|rT|lw(4cm)|rT|lw(4cm)|]
\HL
\VL ~~ER \VL Extra Radical \VL ~~DR \VL Double Radical \VL\SR
\HL
\stoptable

You can comment on a bond. Text is typed in the second argument
of \type{\chemical}. 

\placetable
{Atoms and molecules (radicals).}
\starttable[|rT|lw(4cm)|rT|lw(4cm)|]
\HL
\VL ~~~Z \VL Atom        \VL ~~RZ \VL Radical Atom       \VL\FR
\VL ~CRZ \VL Center Atom \VL ~-RZ \VL Left Radical Atom  \VL\MR
\VL MIDZ \VL Mid Atom    \VL ~+RZ \VL Right Radical Atom \VL\LR
\HL
\stoptable

From these keys \type{RZ} is an addition to the key \type{R}. The key
\type{MID} is only available in combination with a sixring (\type{SIX}). In
the example below we see the effects of \type{MID} and \type{MIDZ}. These
keys have no positioning parameter. 

\startbuffer
\startchemical[frame=on]
  \chemical[SIX,B,MID,MIDZ][\SL{CH_2}]
\stopchemical
\stopbuffer

\startfiguretext
  [left][]
  {none}
  {\getbuffer}
\voorbeeld
\typebuffer
\stopfiguretext

Atoms and molecules are numbered clockwise. Combinations are also allowed.
Position \type{0} (zero) is the middle of a structure.

We can attach labels and numbers to an atom or a bond. This is done with
\type{ZN} and \type{ZT}: 

\placetable
{Labels and numbers.}
\starttable[|rT|lw(4cm)|rT|lw(4cm)|]
\HL
\VL ~~ZN \VL Atom Number \VL ~~ZT \VL Atom Text \VL\SR
\HL
\stoptable

In case of a \type{SIX} and a \type{FIVE} we can also attach text to
radicals. We use \type{RN} and \type{RT}. 

\placetable
{Labels and numbers.}
\starttable[|rT|lw(4cm)|rT|lw(4cm)|]
\HL
\VL ~~RN \VL Radical Number        \VL ~~RT \VL Radical Text        \VL\FR
\VL ~RTN \VL Radical Top Number    \VL ~RTT \VL Radical Top Text    \VL\MR
\VL ~RBN \VL Radical Bottom Number \VL ~RBT \VL Radical Bottom Text \VL\LR
\HL
\stoptable

The structure \type{ONE} has also a top and bottom alternative. 

\placetable
{Extra labels and numbers.}
\starttable[|rT|lw(4cm)|rT|lw(4cm)|]
\HL
\VL ~ZTN \VL Atom Top Number    \VL ~ZTT \VL Atom Top Text    \VL\FR
\VL ~ZBN \VL Atom Bottom Number \VL ~ZBT \VL Atom Bottom Text \VL\LR
\HL
\stoptable

With the keys \type{ZTN} and \type{ZBN} numbers are generated automatically.
The other keys will use the typed text of the second argument of 
\type{chemical}.

\startbuffer
\startchemical[frame=on]
  \chemical[ONE,SB,Z0,ZTT][C,a,b,c,d,e,f,g,h]
\stopchemical
\stopbuffer

\startfiguretext
  [left][]
  {none}
  {\getbuffer}
\voorbeeld
\typebuffer
\stopfiguretext

You can also add some symbols to the structure.

\placetable
{Indications.}
\starttable[|rT|lw(4cm)|rT|lw(4cm)|]
\HL
\VL ~~AU \VL Arrow Up \VL ~~AD \VL Arrow Down \VL\SR
\HL
\stoptable

The arrows are positioned between the atoms in a structure.

\startbuffer
\startchemical[frame=on]
  \chemical[SIX,B,AU]
\stopchemical
\stopbuffer

\startfiguretext
  [left][]
  {none}
  {\getbuffer}
\voorbeeld
\typebuffer
\stopfiguretext

We want to add that while typesetting atoms and molecules the dimensions of
these atoms and molecules are taken into account. The width of \chemical{C}
and the height of \chemical{C^n_m} play an important role during positioning.
This mechanism may be refined in a later stage. 

\chapter{Frontviews}

Structures \type{FIVE} and \type{SIX} can be displayed in a frontview.
However there are some limitations. Frontviews can not be rotated. Also the
coupling of several structures is limited. 

We illustrate the frontview keys in two examples.

\startbuffer
\startchemical[height=4500,bottom=2500]
  \bottext{$\beta$-D-Fructopyranose}
  \chemical
    [SIX,FRONT,BB1236,+SB4,-SB5,Z5][O]
  \chemical
    [SIX,FRONT,+R12346,+RZ12346][\SR{HO},H,H,H,OH]
  \chemical
    [SIX,FRONT,-R12346,-RZ12346][H,OH,\SR{HO},H,CH_2OH]
\stopchemical
\stopbuffer

\startfiguretext
  [left][]
  {none}
  {\getbuffer}
\voorbeeld
\typebuffer
\stopfiguretext

Positioning the radicals is an optimization of feasability and quality. The
next example will illustrate this. 

\startbuffer
\startchemical[height=4500,bottom=2500]
  \bottext{$\alpha$-D-Fructofuranose}
  \chemical
    [FIVE,FRONT,BB125,+SB3,-SB4,Z4][O]
  \chemical
    [FIVE,FRONT,+R1235,+RZ1235][\SR{HO},H,\SR{HOH_2C},CH_2OH]
  \chemical
    [FIVE,FRONT,-R1235,-RZ1235][OH,H,\SR{HO},H,CH_2OH]
\stopchemical
\stopbuffer

\startfiguretext
  [left][]
  {none}
  {\getbuffer}
\voorbeeld
\typebuffer
\stopfiguretext

\chapter{Definitions} 

It is possible to build a library of structures. These predefined structures
can be used in a later stage, for example as a component of a more complex
structure. Predefinition can be done with the \TEX||primitive \type{\def}. 

\starttyping
\def\sixring{\chemical[SIX,B,R,RZ]}
\stoptyping

However it is better to use the command \type{\definechemical}. In that case
a message will occur during processing if a duplicate name is found. 

\starttyping
\definechemical[sixring]
  {\chemical[SIX,B,R,RZ]}
\stoptyping

Recalling \type{\chemical[sixring]} will display a bare sixring without
substituents. 

\startbuffer
\definechemical[sixring]
  {\chemical[SIX,B,R,RZ]}

\startchemical[frame=on,width=6000]
  \chemical[sixring]
\stopchemical
\stopbuffer

\startfiguretext
  [left][]
  {none}
  {\getbuffer}
\voorbeeld
\typebuffer
\stopfiguretext

If we want to attach six substituents in a later stage to a sixring we could
type:

\startbuffer
\definechemical[sixring]
  {\chemical[SIX,B,R,RZ]}

\startchemical[frame=on,width=6000]
  \chemical[sixring][R_1,R_2,R_3,R_4,R_5,R_6]
\stopchemical
\stopbuffer

\startfiguretext
  [left][]
  {none}
  {\getbuffer}
\voorbeeld
\typebuffer
\stopfiguretext

The structure \type{sixring} can be defined without substituents (\type{RZ}).
We could attach them after recalling \type{\chemical[sixring]}. 

\startbuffer
\definechemical[sixring]
  {\chemical[SIX,B,R]}

\startchemical[frame=on,width=6000]
  \chemical[sixring,RZ][A,B,C,D,E,F]
\stopchemical
\stopbuffer

\startfiguretext
  [left][]
  {none}
  {\getbuffer}
\voorbeeld
\typebuffer
\stopfiguretext

In principal the possibilities are unlimited. However, you should remember
that atoms and molecules are selected from the second argument in the order
of definition in the first argument. 

A definition may contain atoms and molecules (texts). 

\starttyping
\definechemical[sixring]
  {\chemical[SIX,B,R,RZ135][R_1,R_3,R_5]}
\stoptyping

In the example above there will always be three substituents. If we want to
attach more substituents we have to indicate explicitly that we want to
continue with the sixring (\type{SIX}). 

\startbuffer
\definechemical[sixring]
  {\chemical[SIX,B,R,RZ135][R_1,R_3,R_5]}

\startchemical[frame=on,width=6000]
  \chemical[sixring,SIX,RZ246][A,B,C]
\stopchemical
\stopbuffer

\startfiguretext
  [left][]
  {none}
  {\getbuffer}
\voorbeeld
\typebuffer
\stopfiguretext

In a definition \type{\chemical[]} has a global scope (this means that
\type{SIX} is remembered) and \type{\chemical[][]} has a local scope. The
idea behind this is that in the first case a range of keys can be added and
in the second case a complete structure. 

In a definition \type{\chemical} may be used more than once. The last example
could have been defined thus: 

\startbuffer
\definechemical[sixring]
  {\chemical[SIX,B,R,RZ135][R_1,R_3,R_5]
   \chemical[SIX,RZ246]}

\startchemical[frame=on,width=6000]
  \chemical[sixring][A,B,C]
\stopchemical
\stopbuffer

\startfiguretext
  [left][]
  {none}
  {\getbuffer}
\voorbeeld
\typebuffer
\stopfiguretext

When \TEX\ announces that an \type{unknown command} has occurred, you may have
forgotten to type \type{SIX}, \type{FIVE} or a comparable key. 

\chapter{Combinations}

Structures can be combined to more complex compounds. Moving one structure in
relation to another structure is done by \type{MOV}, \type{ROT}, \type{ADJ}
and \type{SUB}. 

\placetable
{Moving and rotating.}
\starttable[ | r T | l w(2cm) | p(8cm) | ]
\HL
\VL MOV \VL Move       \VL moving a comparable structure \hbox{structure} in
                           the direction of a bond\VL\FR 
\VL ADJ \VL Adjace     \VL moving another structure in the direction of the
                           $x$-- or $y$--axis, adjacent to a bond\VL\MR 
\VL SUB \VL Substitute \VL moving one structure in relation to another in
                           the direction of the $x$-- or $y$--axis\VL\MR 
\VL ROT \VL Rotate     \VL rotating a structure\VL\LR 
\HL
\stoptable

The four keys mentioned above have different effects when they are applied to
different structures. The angle of rotation in \type{\chemical[FIVE,ROT1,B]}
differs from that in \type{\chemical[SIX,ROT1,B]}. 

With the structure \type{ONE} you can use \type{MOV} but the key \type{DIR}
is also available. Both keys have the same effect but differ in spacing.
Small adjustments are possible with \type{OFF}. 

\placetable
{Moving and rotating.}
\starttable[ | r T | l w(2cm) | p(8cm) | ]
\HL
\VL DIR \VL Direction \VL moving a structure in a diagonal direction\VL\FR
\VL OFF \VL Offset    \VL moving atoms and molecules over small distances\VL\LR 
\HL
\stoptable

The structure \type{CARBON} can be mirrored with \type{MIR}. 

\placetable
{Mirroring.}
\starttable[ | r T | l w(2cm) | p(8cm) | ]
\HL
\VL MIR \VL Mirror     \VL mirroring a structure\VL\SR
\HL
\stoptable

We use a number to indicate the direction of a movement or the level of
rotation. These keys are closely related with the structure. Therefore they
have to be defined before bonds are drawn and texts are placed. So definition
\type{\chemical[FIVE,B,ROT1,R]} and \type{\chemical[FIVE,ROT1,B,R]} will not
have the same result. The first definition will give an undesirable result. 

\startbuffer
\startchemical[frame=on,width=4000,right=3000]
  \chemical[SIX,B,MOV1,B]
\stopchemical
\stopbuffer

\startfiguretext
  [left][]
  {none}
  {\getbuffer}
\voorbeeld
\typebuffer
\stopfiguretext

In this example a sixring is drawn because of \type{SIX,B}. Then a movement
in the direction of bond~1 takes place and a second sixring is
drawn:~\type{B} (\type{SIX} is stil in effect). 

A movement with \type{MOV} in a sixring can occur in six directions. A
movement with \type{ADJ} will take place in only four axis||directions ($x$,
$-x$, $y$, $-y$). It is a coincidence that in a sixring some of these
movements have the same effect. The example above could have been drawn with:
\type{[SIX,B,ADJ1,B]}. 

Structures can be combined. It is possible for example to combine structure
\type{FIVE} with structure \type{SIX} in such a way that they have one mutual
bond. Luckily the mechanism that takes care of these kinds of combinations is
hidden for the user. In the next example you will see a sixring drawn by
\type{SIX,B}. Then a movement in the positive $x$||direction is done by
\type{ADJ1}. At last a rotated fivering is drawn: \type{FIVE,ROT3,B}. 

\startbuffer
\startchemical[frame=on,width=4000,right=3000]
  \chemical[SIX,B,ADJ1,FIVE,ROT3,B]
\stopchemical
\stopbuffer

\startfiguretext
  [left][]
  {none}
  {\getbuffer}
\voorbeeld
\typebuffer
\stopfiguretext

\startbuffer
\startchemical[frame=on,width=5000,right=4500]
  \chemical[SIX,ROT2,B,R6,SUB1,FIVE,B,R4]
\stopchemical
\stopbuffer

\startfiguretext
  [left][]
  {none}
  {\getbuffer}
\voorbeeld
\typebuffer
\stopfiguretext

To go from one structure to an adjacent one is done with \type{ADJ}. Most of
the time one of these structures will have to be rotated to obtain a good
attachment. This is done by \type{ROT}. Rotations are always clockwise in
steps of 90\graden. When a structure is attached with a bond you will have to
use \type{SUB}. Movements with \type{ADJ} and \type{SUB} take place in the
four directions of the $x$- and $y$||axis. 

The next examples illustrate that the dimensions of the smaller structures
are determined by the larger structures, especially \type{SIX}. You will
notice that \type{EIGHT} has fewer possibilities than \type{SIX}. 

\startbuffer
\startchemical[width=6000,left=1500,frame=on]
  \chemical[EIGHT,B,MOV1,B]
\stopchemical
\stopbuffer

\startfiguretext
  [left][]
  {none}
  {\getbuffer}
\voorbeeld
\typebuffer
\stopfiguretext

\startbuffer
\startchemical[width=6000,left=1500,frame=on]
  \chemical[EIGHT,B,ADJ1,SIX,B]
\stopchemical
\stopbuffer

\startfiguretext
  [left][]
  {none}
  {\getbuffer}
\voorbeeld
\typebuffer
\stopfiguretext

\startbuffer
\startchemical[width=6000,left=1500,frame=on]
  \chemical[EIGHT,B,ADJ1,FIVE,ROT3,B]
\stopchemical
\stopbuffer

\startfiguretext
  [left][]
  {none}
  {\getbuffer}
\voorbeeld
\typebuffer
\stopfiguretext

It will be clear by now that the order in which the keys are defined 
makes a lot of difference. The order should be:

\starttyping
\chemical
  [structure,                              % SIX, FIVE, ...
   bonds within the structure,             % B, C, EB, ...
   bonds pointing to substituents,         % R, DR, ...
   atoms within the structure,             % Z
   subsituents attached to the structure]  % RZ, -RZ, ...
  [atoms,
   substituents]
\stoptyping

Most of the time putting structures together is done by translating and
rotating. You could automate this process. In earlier versions this was done
automatically, however this led to misinterpretations of users concerning the
positions of bonds, atoms and substituents. A structure that consists of more
than one component can best be defined per component, translations included.
Rotations should wait until the last step. 

A sixring may have substituents consisting of a carbon chain. In those 
situations we use \type{DIR} to build the chain. 

\startbuffer
\startchemical
  [scale=small,width=6000,height=6000,frame=on]
  \chemical[SIX,SB2356,DB14,Z2346,SR36,RZ36]
    [C,N,C,C,H,H_2]
  \chemical[PB:Z1,ONE,Z0,DIR8,Z0,SB24,DB7,Z27,PE]
    [C,C,CH_3,O]
  \chemical[PB:Z5,ONE,Z0,DIR6,Z0,SB24,DB7,Z47,PE]
    [C,C,H_3C,O]
  \chemical[SR24,RZ24]
    [CH_3,H_3C]
\stopchemical
\stopbuffer

\startfiguretext
  [left][]
  {none}
  {\getbuffer}
\voorbeeld
\typebuffer
\stopfiguretext

Because chains have no predefined format the chains are build and positioned
as a substructure. For positioning we use the keys \type{PB} and \type{PE}. 

\placetable
{Positioning.}
\starttable[ | r T | l w(4cm) | p(6cm) | ]
\HL
\VL PB:.. \VL Picture Begin \VL beginning a substructure \VL\FR
\VL PE    \VL Picture End   \VL ending a substructure    \VL\LR
\HL
\stoptable

Directly after \type{PB} you will have to define the location where the
substructure is positioned. The first following atom is centered on that
location. Always use a central atom on this location. 

These keys were introduced after trying to obtain structures that are
typeset in an acceptable quality. There are different ways to define
structures. The following alternative would have resulted in: 

\startbuffer
\startchemical
    [scale=small,width=6000,height=6000,frame=on]
  \chemical
    [SIX,SB2356,DB14,Z36,SR36,RZ36][N,C,H,H_2]
  \chemical
    [PB:Z1,ONE,Z0,DIR8,Z0,SB24,DB7,Z27,PE][C,C,CH_3,O]
  \chemical
    [PB:Z5,ONE,Z0,DIR6,Z0,SB24,DB7,Z47,PE][C,C,H_3C,O]
  \chemical
    [PB:Z2,ONE,Z0,DIR2,SB6,CZ0,PE][C,CH_3]
  \chemical
    [PB:Z4,ONE,Z0,DIR4,SB8,CZ0,PE][C,H_3C]
\stopchemical
\stopbuffer

\startfiguretext
  [left][]
  {none}
  {\getbuffer}
\voorbeeld
\typebuffer
\stopfiguretext


The most efficient way to define such a structure would be like the example
below. Typographically you wouldn't be satisfied. 

\startbuffer
\startchemical
    [scale=small,width=6000,height=6000,frame=on]
  \chemical[SIX,SB2356,DB14,Z,SR36,RZ36,SR1245,RZ24]
    [C,C,N,C,C,C,H,H_2,CH_3,H_3C]
  \chemical[PB:RZ1,ONE,Z0,SB2,DB7,Z27,PE]
    [C,CH_3,O]
  \chemical[PB:RZ5,ONE,Z0,SB4,DB7,Z47,PE]
    [C,H_3C,O]
\stopchemical
\stopbuffer

\startfiguretext
  [left][]
  {none}
  {\getbuffer}
\voorbeeld
\typebuffer
\stopfiguretext

You may have noticed that the measurements of the structure is determined by
the substituents. The chains are not taken into account. This leads to a
consistent build-up of a structure. 

The differences in outcome when using \type{SUB} in stead of \type{PB} are
very small. However compare the following formulas. 

\startbuffer
\startchemical
  [width=fit,frame=on,scale=small]
  \chemical
    [SIX,ROT2,B,C,R236,RZ23,
     SUB1,ONE,OFF1,Z0,4OFF1,SB1,Z1]
    [HO,HO,CHOH,CH_2NH_2]
\stopchemical
\stopbuffer

\startfiguretext
  [left][]
  {none}
  {\getbuffer}
\voorbeeld
\typebuffer
\stopfiguretext

\startbuffer
\startchemical
  [width=fit,frame=on,scale=small]
  \chemical
    [SIX,ROT2,B,C,R236,RZ23,
     PB:RZ6,ONE,Z0,3OFF1,SB1,Z1,PE]
    [HO,HO,CHOH,CH_2NH_2]
\stopchemical
\stopbuffer

\startfiguretext
  [left][]
  {none}
  {\getbuffer}
\voorbeeld
\typebuffer
\stopfiguretext

The use of the key \type{PB:} might be somewhat more difficult, but the
results are much better. In that case you should define the width yourself,
because the substituents are not taken into account when determining the
dimensions. 

First we will go into the key \type{OFF}. In some cases atom (\type{Z0}) in
\type{ONE} can consist of more than one character. The reserved space for
these characters would be insufficient and character and bond would overlap.
When you need more space for \type{Z0} we can move bond~1, 2 and~8 by means
of the key~\type{OFF} ('offset'). The example below will illustrate its use. 

\startbuffer
\startchemical
  [width=fit,size=small,scale=small,frame=on]
  \chemical
    [SIX,B,C,ADJ1,
     FIVE,ROT3,SB34,+SB2,-SB5,Z345,DR35,SR4,CRZ35,SUB1,
     ONE,OFF1,SB258,Z0,Z28]
    [C,N,C,O,O,
     CH,COOC_2H_5,COOC_2H_5]
\stopchemical
\stopbuffer

\startfiguretext
  [left][]
  {none}
  {\getbuffer}
\voorbeeld
\typebuffer
\stopfiguretext

Moving the bonds makes room for an extra character. More space was obtained
when we would have typed \type{3OFF1}. The example looks rather complex but
you can define it rather easy by defining its components first. Rotating
should be done in the last stage. 

You see a new key: \type{CRZ}. This key is used to place the atom or molecule
in one line with the bond. You could have used \type{RZ}, because you can
influence spacing in the second argument with \type{{\,O}} in stead off
\type{O} (spacing in mathematical mode). 

We will show another example, produced in two ways. When choosing a method
you should take into account the consistency throughout your document. 

\startbuffer
\startchemical
  [width=fit,height=fit,frame=on,
   scale=small]
  \chemical
    [ONE,SB15,DB7,Z057,3OFF1,MOV1,Z0,3OFF1,MOV1,
     Z017,SB1357,MOV3,Z0,MOV3,SB1357,Z013,3OFF5,
     MOV5,Z0,3OFF5,SB5,Z5]
    [C,H_2N,NH,(CH_2)_3,C,COOH,H,\SL{NH},C,COOH,H,
     (CH_2)_2,HOOC]
\stopchemical
\stopbuffer

\startfiguretext
  [left][]
  {none}
  {\getbuffer}
\voorbeeld
\typebuffer
\stopfiguretext

\startbuffer
\startchemical
  [width=fit,height=fit,frame=on,
   scale=small]
  \chemical [ONE,SB15,DB7,Z057,3OFF1][C,H_2N,NH]
  \chemical [MOV1,Z0,3OFF1]          [(CH_2)_3]
  \chemical [MOV1,Z017,SB1357]       [C,COOH,H]
  \chemical [MOV3,Z0]                [\SL{NH}]
  \chemical [MOV3,SB1357,Z013,3OFF5] [C,COOH,H]
  \chemical [MOV5,Z0,3OFF5,SB5,Z5]   [(CH_2)_2,HOOC]
\stopchemical
\stopbuffer

\startfiguretext
  [left][]
  {none}
  {\getbuffer}
\voorbeeld
\typebuffer
\stopfiguretext

\startbuffer
\startchemical
  [width=fit,height=fit,frame=on,
   scale=small]
  \chemical
    [ONE,Z0,SAVE,MOV7,SB1357,Z017,3OFF5,MOV5,Z0,
     3OFF5,MOV5,SB15,DB7,Z057,RESTORE,
     MOV3,SB1357,Z013,MOV5,3OFF5,Z0,6OFF5,SB5,Z5]
    [\SL{NH},C,COOH,H,(CH_2)_3,C,H_2N,NH,C,COOH,H,
     (CH_2)_2,HOOC]
\stopchemical
\stopbuffer

\startfiguretext
  [left][]
  {none}
  {\getbuffer}
\voorbeeld
\typebuffer
\stopfiguretext

\startbuffer
\startchemical
  [width=fit,height=fit,frame=on,scale=small]
  \chemical
    [ONE,Z0,MOV7,SB1357,Z017,3OFF5,MOV5,Z0,
     3OFF5,MOV5,SB15,DB7,Z057,MOV0,MOV3,SB1357,
     Z013,MOV5,3OFF5,Z0,6OFF5,SB5,Z5]
    [\SL{NH},C,COOH,H,(CH_2)_3,C,H_2H,NH,C,COOH,H,
     (CH_2)_2,HOOC]
\stopchemical
\stopbuffer

\startfiguretext
  [left][]
  {none}
  {\getbuffer}
\voorbeeld
\typebuffer
\stopfiguretext

Notice the use of \type{SAVE} and \type{RESTORE}. These keys enable you to
save a location in a structure and return to that location in another stage. 

As an extra we will show you a combination of \type{SIX} and \type{FIVE}. Be
aware of the use of \type{SS}. 

\startbuffer
\startchemical
    [width=fit,height=fit,frame=on]
  \chemical
    [SIX,DB135,SB246,Z,SR6,RZ6][C,C,N,\SR{HC},N,C,NH_2]
  \chemical
    [SIX,MOV1,DB1,SB23,SS6,Z1..5,SR3,RZ3][N,\SL{CH},N,C,C,H]
\stopchemical
\stopbuffer

\startfiguretext
  [left][]
  {none}
  {\getbuffer}
\voorbeeld
\typebuffer
\stopfiguretext

\chapter{Extra text}

We can add text and symbols in and around structures. For example:

\startbuffer
\startchemical
  [height=4500,top=1250,width=fit,frame=on]
  \bottext
    {2,2-Dimethyl-3-ethylpentane}
  \chemical
    [ONE,Z3570,SB1357]
    [CH_3,\T{1}{H_3C},CH_3,\SR{\LT{2}{C}}]
  \chemical
    [MOV1,OFF1,Z0,SB3]
    [\T{3}{CH}]
  \chemical
    [MOV3,Z0,SB3,MOV3,Z0,MOV7,MOV7]
    [CH_2,CH_3]
  \chemical
    [OFF1,SB1,MOV1,OFF1,Z0,2OFF1,SB1,Z1]
    [\T{4}{CH_2},\T{5}{CH_3}]
\stopchemical
\stopbuffer

\startfiguretext
  [left][]
  {none}
  {\getbuffer}
\voorbeeld
\typebuffer
\stopfiguretext

There is a range of keys like \type{\T}. In a number of cases the arguments
are optional. Charges can be displayed in Roman by means of \type{\+} and
\type{\-} or directly by means of~\type{\1} up to~\type{\7}. 

\placetable
{Text: charges.}
\starttable[ | l | l | ]
\HL
\VL \type{\+{number}} \VL positive charge in Roman   \VL\FR
\VL \type{\-{number}} \VL negative charge in Roman   \VL\MR
\VL \type{\1}         \VL I (without sign)           \VL\MR
\VL \type{\7}         \VL II, III, IV, V, VI and VII \VL\LR
\HL
\stoptable

A charge is centered above the atom. For example:

\startbuffer
\placeformula
  \startformula
    \chemical{S}+\chemical{O_2}
    \chemical{GIVES}
    \chemical{\+{4}{S}\-{2}{O_2}}
  \stopformula
\stopbuffer

\typebuffer

will result in:

\getbuffer

If we want to repeat a number of atoms or molecules we can define an
(endless) range with \type{\[} and \type{\]}. Both arguments are optional as
is shown in the example of \kap{PTFE} of Polytetrafluorethane, better knowns
as Teflon. 

\startbuffer
\startchemical[width=fit,frame=on]
  \chemical
    [ONE,ZT5,SB5,OFF1,Z0,OFF1,SB1,MOV1,SB5,OFF1,Z0,OFF1,SB1,ZT1]
    [\[,CF_2,CF_2,\]{\sl n}]
\stopchemical
\stopbuffer

\startfiguretext
  [left][]
  {none}
  {\getbuffer}
\voorbeeld
\typebuffer
\stopfiguretext

\placetable
{Text: repeating.}
\starttable[ | l | l | l | ]
\HL
\VL \type{\[{bottom}} \VL \type{\[{top}{bottom}} \VL right repeating sign \VL\FR
\VL \type{\]{bottom}} \VL \type{\]{top}{bottom}} \VL left repeating sign  \VL\LR
\HL
\stoptable


There is no problem of placing texts 
on the left, right, top or bottom of the atoms or molecules.
If we preceed the keys \type{\L}, \type{\R}, \type{\T} and \type{\B} 
by \type{\X} the distance from text to atoms is somewhat smaller.

\placetable
{Text: around an atom.}
\starttable[ | l | l | ]
\HL
\VL \type{\L{text}} \VL text left   \VL\FR
\VL \type{\R{text}} \VL text right  \VL\MR
\VL \type{\T{text}} \VL text top    \VL\MR
\VL \type{\B{text}} \VL text bottom \VL\LR
\HL
\stoptable

Logical combinations of these keys are also possible. A key to centre text is
also available. 

\leavevmode\hbox to \hsize
  {\def\sample#1{\chemical{#1{\oplus}{\ruledhbox{\ttx\string#1}}}}%
   \sample\TL\hss\sample\L \hss\sample\LC\hss\sample\BL\hss
   \sample\TR\hss\sample\R \hss\sample\RC\hss\sample\BR\hss
   \sample\LT\hss\sample\T \hss\sample\RT\hss\sample\LB\hss
   \sample\B \hss\sample\RB}

\leavevmode\hbox to \hsize
  {\def\sample#1{\chemical{\X#1{\oplus}{\ruledhbox{\ttx\string\X\string#1}}}}%
   \sample\TL\hss\sample\L \hss\sample\LC\hss\sample\BL\hss
   \sample\TR\hss\sample\R \hss\sample\RC\hss\sample\BR\hss
   \sample\LT\hss\sample\T \hss\sample\RT\hss\sample\LB\hss
   \sample\B \hss\sample\RB}

In some cases you will need what we may call {\em smashed} text. 

\placetable
{Text: smashed text.}
\starttable[ | l | l | ]
\HL
\VL \type{\SL{text}} \VL left align  \VL\FR
\VL \type{\SM{text}} \VL centre      \VL\MR
\VL \type{\SR{text}} \VL right align \VL\LR
\HL
\stoptable

An example is given below. The text is centred around the first character.

\startbuffer
\startchemical[frame=on]
  \chemical[SIX,B,R36,RZ36][\SL{COOH},\SL{COOH}]
\stopchemical
\stopbuffer

\startfiguretext
  [left][]
  {none}
  {\getbuffer}
\voorbeeld
\typebuffer
\stopfiguretext

We can place text above, under or in the middle of
structures with the keys \type{\toptext}, \type{\midtext}
and \type{\bottext}. The exact position is determined by the
height and depth of the structure. 

\startbuffer
\placeformula
  \startformula
    \startchemical[width=fit]
      \chemical[SIX,B,Z1,MOV1,B][\hbox{$\bullet$}]
      \toptext{{\sl trans}-Decalin}
    \stopchemical
    \hskip 24pt
    \startchemical[width=fit]
      \chemical[SIX,B,Z12,MOV1,B][\hbox{$\bullet$},\hbox{$\bullet$}]
      \bottext{{\sl cis}-Decalin}
    \stopchemical
  \stopformula
\stopbuffer

\typebuffer

Both {\em Decalin} formulas look like this:

\getbuffer 

\chapter{Axis}

Structures can be typeset in a frame that is divided by axis. The dimensions
of the axis and the location of the origin can be defined in the set up. The
axis and the frame can be made visible. 

\startbuffer
\startchemical
  [axis=on,
   width=6000,height=4000]
\stopchemical
\stopbuffer

\startfiguretext
  [left][]{none}
  {\getbuffer}
\voorbeeld[ex:axis]
\typebuffer
\stopfiguretext

\startbuffer
\startchemical
  [axis=on,
   left=2000,right=4000]
\stopchemical
\stopbuffer

\startfiguretext
  [left][]{none}
  {\getbuffer}
\voorbeeld
\typebuffer
\stopfiguretext

\startbuffer
\startchemical
  [axis=on,
   width=6000,right=1000]
\stopchemical
\stopbuffer

\startfiguretext
  [left][]{none}
  {\getbuffer}
\voorbeeld
\typebuffer
\stopfiguretext

\startbuffer
\startchemical
  [axis=on,
   width=6000,top=1000,bottom=3000]
\stopchemical
\stopbuffer

\startfiguretext
  [left][]{none}
  {\getbuffer}
\voorbeeld
\typebuffer
\stopfiguretext

The dimensions of the total structure determine the dimensions of the axis.
When \type{width=fit} and/or \type{height=fit}
is typed the dimensions are determined by the real dimensions. Your choice 
will depend on how you want to place the structure in the text.

\in{Example}[ex:axis] shows the default set up. Within a
\type{\start}||\type{\stop}||pair you can use \PICTEX||macros. However, be
careful. 

\chapter{Set ups}

After \type{\startchemical} and  \type{\setupchemical} you can type the 
set up.

\placetable
{Set ups for structures.}
\starttable[|lT|Tl|lT|]
\HL
\VL \bf parameter \VL \bf values            \VL \bf default \VL\SR
\HL
\VL width         \VL number                     \VL 4000   \VL\FR
\VL height        \VL number                     \VL 4000   \VL\MR
\VL left          \VL number                     \VL        \VL\MR
\VL right         \VL number                     \VL        \VL\MR
\VL top           \VL number                     \VL        \VL\MR
\VL bottom        \VL number                     \VL        \VL\LR
\HL
\VL resolution    \VL number                     \VL \type{\outputresolution} \VL\FR
\VL bodyfont      \VL 8pt 9pt 10pt etc.          \VL \type{\bodyfontsize}        \VL\MR
\VL character     \VL \type{\rm} \type{\bf} etc. \VL \type{\rm}               \VL\LR
\HL
\VL scale         \VL number                     \VL medium \VL\FR
\VL size          \VL small medium big           \VL medium \VL\LR
\HL
\VL state         \VL start stop                 \VL start  \VL\FR
\VL option        \VL test                       \VL        \VL\LR
\HL
\VL axis          \VL on off                     \VL off    \VL\FR
\VL frame         \VL on off                     \VL off    \VL\LR
\HL
\VL alternative   \VL 1 2                        \VL 1      \VL\SR
\HL
\VL offset        \VL HIGH LOW                   \VL LOW    \VL\SR
\HL
\stoptable

The axis range from $-2000$ upto $+2000$, height as well as width. The
parameter \type{Z0} is at $(0,0)$. Other divisions can be set up with
\type{left}, \type{right}, \type{top} and/or \type{bottom} in combination
with \type{width} and \type{height}. 

You can use the key \type{size} to set up the bodyfont. In doing so the
\TEX||primitives \type{\textsize}, \type{\scriptsize} and
\type{\scriptscriptsize} are used. With \type{scale} you can set up the
dimensions of the structure itself. The scale is determined by the parameter
\type{bodyfont}. The values \type{small}, \type{medium} and \type{big} are
proportionally related. 

The set up of the parameter \type{bodyfont} is used for calculations and has
no consequences for the text. In \CONTEXT\ and \LATEX\ this set up parameter
is coupled to the mechanism that sets the bodyfont. 

In \TEX\ and \CONTEXT\ you can use commands like \type{\rm}, \type{\bf} and
\type{\sl} in mathematical mode. \PPCHTEX\ uses default \type{\rm}. With the
parameter \type{character} another alternative can be chosen. In
\in{example}[ex:character] the substituents are typeset {\sl slanted}. 

\startbuffer
\startchemical[frame=on,character=\sl]
  \chemical[SIX,B,R,RZ][A,B,C,D,E,F]
\stopchemical
\stopbuffer

\startfiguretext
  [left][]
  {none}
  {\getbuffer}
\voorbeeld[ex:character]
\typebuffer
\stopfiguretext

The set up of \type{character} is valid for chemical structures in a picture
and in the text. The sub- and superscripts are changed accordingly. This is
illustrated in {\setupchemical[character=\rm]\chemical{CH_4}},
{\setupchemical[character=\bf]\chemical{CH_4}} and
{\setupchemical[character=\sl]\chemical{CH_4}}, in which the set ups are
\type{\rm}, \type{\bf} and \type{\sl}. Italic \type{\it} formulas lead to a
bigger linewidth. In \CONTEXT\ default bold||slanted (\type{\bs}) and
bold||italic (\type{\bi}) are available. These commands adjust automatically
to the actual fontstyle: {\setupchemical[character=\ss\sl]\chemical{CH_4}},
{\setupchemical[character=\rm\sl]\chemical{CH_4}},
{\setupchemical[character=\tt\sl]\chemical{CH_4}} etc. (\type{\ss}, \type{\rm},
\type{\tt}). 

It is also possible to set the characters at the instant you type them in 
the argument.

\startbuffer
\startchemical[frame=on]
  \chemical[SIX,B,R,RZ][\tf a,\bf b,\it c,\sl d,\bi e,\bs f]
\stopchemical
\stopbuffer

\startfiguretext
  [left][]
  {none}
  {\getbuffer}
\voorbeeld
\typebuffer
\stopfiguretext

With parameter \type{state} calculations can be shortcut. The parameters
\type{frame} and \type{axis} need no further explanation. With
\type{option=test} frames are drawn around the texts in a structure. In this
way you can see how the text is aligned.\footnote{In \CONTEXT\ you can
activate the visual debugger. When activated the baseline is a dotted line.
The module \type{supp-vis} can be used in other systems.} With the parameter
\type{alternative} you can set up the quality of the lines. Default \PICTEX\ uses
a 5~point . to draw lines. When option~2 is chosed a thinner line is drawn. 

\startbuffer
\startchemical[frame=on,option=test,alternative=2]
  \chemical[SIX,B,R,RZ][R_1,R_2,R_3,R_4,R_5,R_6]
\stopchemical
\stopbuffer

\startfiguretext
  [left][]
  {none}
  {\getbuffer}
\voorbeeld[ex:test]
\typebuffer
\stopfiguretext

The offset relates to the position of the sub- and superscripts. With
\type{HIGH} the subscripts are placed high (\chemical{HIGH,H_2O}) and with
\type{LOW} somewhat lower (\chemical{LOW,H_2O}). 

A structure can be displayed in different sizes. This is done with
\type{formaat} and \type{scale}. 

\startbuffer
\startchemical[frame=on,scale=small,size=small]
  \chemical[SIX,B,R,RZ][1,2,3,4,5,6]
\stopchemical
\stopbuffer

\startfiguretext
  [left][]
  {none}
  {\getbuffer}
\voorbeeld
\typebuffer
\stopfiguretext

\startbuffer
\startchemical[frame=on,scale=medium,size=medium]
  \chemical[SIX,B,R,RZ][1,2,3,4,5,6]
\stopchemical
\stopbuffer

\startfiguretext
  [left][]
  {none}
  {\getbuffer}
\voorbeeld
\typebuffer
\stopfiguretext

\startbuffer
\startchemical[frame=on,scale=big,size=big]
  \chemical[SIX,B,R,RZ][1,2,3,4,5,6]
\stopchemical
\stopbuffer

\startfiguretext
  [left][]
  {none}
  {\getbuffer}
\voorbeeld
\typebuffer
\stopfiguretext

You can also type a number between~1 and 1000 in \type{scale}. The values
belonging to \type{small}, \type{medium} or \type{big} are proportionally 
related.

\chapter{Symbols}

There are some symbols that can be used to display reactions. The reaction
below is typed by: 

\startbuffer
\setupchemical
  [size=small,
   scale=small,
   width=fit,
   height=5500,
   bottom=1500]

\hbox
  {\startchemical
     \chemical[SIX,B,ER6,RZ6][O]
   \stopchemical
   \startchemical
     \chemical[SPACE,PLUS,SPACE]
   \stopchemical
   \startchemical
     \chemical[FIVE,ROT4,B125,+SB3,-SB4,Z4,SR4,RZ4][N,H]
   \stopchemical
   \startchemical
     \chemical[SPACE,GIVES,SPACE][?]
   \stopchemical
   \startchemical
     \chemical[SIX,B,EB6,R6,SUB4,FIVE,ROT4,B125,+SB3,-SB4,Z4][N]
   \stopchemical
   \startchemical
     \chemical[SPACE,PLUS,SPACE,CHEM][H_2O]
   \stopchemical}
\stopbuffer

\typebuffer

The \type{\hbox} is necessary to align the structures. The symbols
\type{GIVES} and \type{PLUS} need no further explanation. With \type{SPACE}
more room can be created between the structures and symbols. 

\placefigure
  [hier]
  [fig:reactie]
  {none}
  {\getbuffer}

An equilibrium can be displayed with \type{EQUILIBRIUM}. Over \type{GIVES}
and \type{EQUILIBRIUM} you can place text. In the example the text is just a
'?'. In addition \type{MESOMERIC} is also available. Braces used for
displaying complexes can be created with \type{OPENCOMPLEX} and
\type{CLOSECOMPLEX}. 

\chapter{Positioning}

When you are combining atoms or molecules, for example with \type{SUB}, some
positions and dimensions change their value. To overcome this problem it is
possible to save a location with \type{SAVE} and return to that location with
\type{RESTORE}. 

\placetable
{Positioning.}
\starttable[ | r T | l w(4cm) | p(6cm) | ]
\HL
\VL SAVE    \VL Save Status    \VL save actual status    \VL\FR
\VL RESTORE \VL Restore Status \VL restore actual status \VL\LR
\HL
\stoptable

The keys \type{SAVE} and \type{RESTORE} are used with substituents. When
placing radicals we use \type{PB} and \type{PE}. This example also
illustrates the possibility to create chains. 

\startbuffer
\definechemical[molecule]
  {\chemical
     [ONE,Z0,SB1357,
      SAVE,SUB2,SIX,B,R6,C,RESTORE,
      MOV1,Z0,SB137,
      MOV1,Z0,SB37,
      MOV1]
     [C,C,C]}

\startchemical[width=fit,height=fit]
  \chemical[molecule,molecule,molecule]
\stopchemical
\stopbuffer

\typebuffer

\placefigure
  [here,force]
  [fig:save]
  {none}
  {\getbuffer}

The example below is more complicated and show a complete reaction. The set 
up of bottom and top is essential in this example.

\startbuffer
\placeformula
  \startformula
    \setupchemical
      [width=fit,top=2000,bottom=2000,
       scale=small,size=small]
    \startchemical
      \chemical
        [ONE,
         SAVE,
           Z0,SB7,SB3,SB1,MOV1,Z0,SB1,MOV1,Z0,DB8,CZ8,SB1,Z1,
         RESTORE,
         SAVE,
           SUB4,ONE,Z0,SB3,SB1,MOV1,Z0,SB1,MOV1,Z0,DB8,CZ8,SB1,Z1,
         RESTORE,
         SUB2,ONE,Z0,SB7,SB1,MOV1,Z0,SB1,MOV1,Z0,DB8,CZ8,SB1,Z1]
        [\SR{HC},O,C,O,C_{19}H_{39},
         \SR{H_{2}C},O,C,O,C_{17}H_{29},
         \SR{H_{2}C},O,C,O,C_{21}H_{41}]
    \stopchemical
    \startchemical
      \chemical[SPACE,PLUS,SPACE]
    \stopchemical
    \startchemical[right=600]
      \chemical[ONE,CZ0][3CH_{3}OH]
    \stopchemical
    \startchemical
      \chemical[SPACE,GIVES,SPACE,SPACE][H^+/H_2O]
    \stopchemical
    \startchemical
      \chemical
        [ONE,
         SAVE,
           Z0,SB7,SB3,SB1,Z1,
         RESTORE,
         SAVE,
           SUB4,ONE,Z0,SB3,SB1,Z1,
         RESTORE,
         SUB2,ONE,Z0,SB7,SB1,Z1]
        [\SR{HC},OH,
         \SR{H_{2}C},OH,
         \SR{H_{2}C},OH]
    \stopchemical
    \startchemical
      \chemical[SPACE,PLUS,SPACE]
    \stopchemical
    \startchemical
      \chemical
        [ONE,
         SAVE,
           Z0,DB8,CZ8,SB1,SB5,Z5,MOV1,Z0,SB1,Z1,
         RESTORE,
         SAVE,
           SUB4,ONE,Z0,DB8,CZ8,SB1,SB5,Z5,MOV1,Z0,SB1,Z1,
         RESTORE,
         SUB2,ONE,Z0,DB8,CZ8,SB1,SB5,Z5,MOV1,Z0,SB1,Z1]
        [C,O,C_{19}H_{39},O,CH_{3},
         C,O,C_{17}H_{29},O,CH_{3},
         C,O,C_{21}H_{41},O,CH_{3}]
    \stopchemical
  \stopformula
\stopbuffer

\typebuffer

This definition might have been more compact if we would have typed
\type{SB731} in stead of \type{SB7,SB3,SB1}. But in this way the definition
is readable. Complex structures can best be defined in its respective
components. 

\getbuffer

Just two more examples where we place text under a structure. 

\startbuffer
\placeformula
  \startformula
    \setupchemical
      [width=fit,top=1500,bottom=3500]
    \startchemical
      \chemical
        [ONE,Z0,DB1,SB3,SB7,Z7,MOV1,Z0,SB3,SB7,Z3,Z7,
         MOV0,SUB2,SIX,B,R6,C]
        [C,H,C,H,H]
      \bottext{styreen}
    \stopchemical
    \quad\quad\quad
    \startchemical
      \chemical
        [ONE,Z0,DB1,SB3,SB7,Z3,Z7,
         MOV1,Z0,SB1,SB3,Z3,
         MOV1,Z0,DB1,SB3,Z3,
         MOV1,Z0,SB3,SB7,Z3,Z7]
        [C,H,H,C,H,C,H,C,H,H]
      \bottext{1,3-butadieen}
   \stopchemical
  \stopformula
\stopbuffer

\typebuffer

\getbuffer

The use of \type{OFF} can be very subtle. The examples below illustrate this
and show minor shifts of \type{ONE}. 

\startbuffer
\startchemical[height=fit,frame=on]
  \chemical[ONE,SB]
\stopchemical
\stopbuffer

\startfiguretext
  [left][]
  {none}
  {\getbuffer}
\voorbeeld
\typebuffer
\stopfiguretext

\startbuffer
\startchemical[height=fit,frame=on]
  \chemical[ONE,3OFF1,SB]
\stopchemical
\stopbuffer

\startfiguretext
  [left][]
  {none}
  {\getbuffer}
\voorbeeld
\typebuffer
\stopfiguretext

\startbuffer
\startchemical[height=fit,frame=on]
  \chemical[ONE,MOV1,SB]
\stopchemical
\stopbuffer

\startfiguretext
  [left][]
  {none}
  {\getbuffer}
\voorbeeld
\typebuffer
\stopfiguretext

\startbuffer
\startchemical[height=fit,frame=on]
  \chemical[ONE,3OFF1,MOV1,SB]
\stopchemical
\stopbuffer

\startfiguretext
  [left][]
  {none}
  {\getbuffer}
\voorbeeld
\typebuffer
\stopfiguretext

\startbuffer
\startchemical[height=fit,frame=on]
  \chemical[ONE,MOV1,3OFF1,SB]
\stopchemical
\stopbuffer

\startfiguretext
  [left][]
  {none}
  {\getbuffer}
\voorbeeld
\typebuffer
\stopfiguretext

\startbuffer
\startchemical[height=fit,frame=on]
  \chemical[ONE,MOV1,3OFF1,OFF0,SB]
\stopchemical
\stopbuffer

\startfiguretext
  [left][]
  {none}
  {\getbuffer}
\voorbeeld
\typebuffer
\stopfiguretext

\startbuffer
\startchemical[height=fit,frame=on]
  \chemical[ONE,MOV1,3OFF1,MOV0,SB] 
\stopchemical
\stopbuffer

\startfiguretext
  [left][]
  {none}
  {\getbuffer}
\voorbeeld
\typebuffer
\stopfiguretext

\startbuffer
\startchemical[height=fit,frame=on]
  \chemical[ONE,MOV1,MOV0,SB]
\stopchemical
\stopbuffer

\startfiguretext
  [left][]
  {none}
  {\getbuffer}
\voorbeeld
\typebuffer
\stopfiguretext

The next example shows the definition of complexes. Pay special attention to
the use of \type{RBT}. Normally an extra spacing is not necessary but we use
here |<|the command is not visible|>| a smaller bodyfont to prevent the
structure to run in the margin. 

\startbuffer
\startformula
\setupchemical[scale=small,width=fit]
\startchemical
  \chemical[OPENCOMPLEX,SPACE]
\stopchemical
\startchemical
  \chemical[SIX,B,EB35,R6,-R1,+R1]                  
  \chemical[SIX,PB:RZ6,ONE,OFF1,Z0,EP57,PE][\SL{OH}]
  \chemical[SIX,-RZ1,+RZ1,RT2][H,Br,\oplus]
\stopchemical
\startchemical
  \chemical[SPACE,MESOMERIC,SPACE]
\stopchemical
\startchemical
  \chemical[SIX,B,EB25,R6,-R1,+R1]                  
  \chemical[SIX,PB:RZ6,ONE,OFF1,Z0,EP57,PE][\SL{OH}] 
  \chemical[SIX,-RZ1,+RZ1,RT4][H,Br,\oplus]
\stopchemical
\startchemical
  \chemical[SPACE,MESOMERIC,SPACE]
\stopchemical
\startchemical
  \chemical[SIX,B,EB24,R6,-R1,+R1]
  \chemical[SIX,PB:RZ6,ONE,OFF1,Z0,EP57,PE][\SL{OH}]  
  \chemical[SIX,-RZ1,+RZ1,RBT6][H,Br,\ \oplus]
\stopchemical
\startchemical
  \chemical[SPACE,MESOMERIC,SPACE]
\stopchemical
\startchemical
  \chemical[SIX,B,EB24,ER6,-R1,+R1]                   
  \chemical[SIX,PB:RZ6,ONE,OFF1,Z0,EP5,ZT7,PE][\SL{OH},\oplus]
  \chemical[SIX,-RZ1,+RZ1][H,Br]
\stopchemical
\startchemical
  \chemical[SPACE,CLOSECOMPLEX]
\stopchemical
\stopformula
\stopbuffer

\typebuffer

Without the use of \type{SPACE} the seperate structures would merge. Most of
the time the optimization of such a reaction is an iterative proces. 

\start
\switchtobodyfont[small]
\getbuffer
\stop

\chapter{Reactions}

Not only the typesetting of chemical structures is supported but also the
typesetting of normal reactions. The command \type{\chemical} has three other
appearances: 

\starttyping
\chemical{formula}
\chemical{formula}{bottom text}
\chemical{formula}{top text}{bottom text}
\stoptyping

This command adapts itself to text mode. That means that it
'knows' whether it is used in: 

\startitemize[packed]
\item text||mode
\item mathematical text||mode
\item mathematical display||mode
\stopitemize

When the command is used in running text it will automatically be surrounded
by \type{$ $}. Typing \type{\chemical{NH_4^+}} will result in
\chemical{NH_4^+}. 

The result would be the same if we would place the command between \type{$
$}. In both cases the second and third argument can be left out. If we place
the command between \type{$$ $$} (or \type{\startformula} and
\type{\stopformula}) both arguments do have a function. First a simple
example. The command \type{\placeformula} is a \CONTEXT\ command and handles
the positioning and numbering of the formula. 

\startbuffer
\placeformula
  \startformula
    \chemical{2H_2} \chemical{PLUS} \chemical{O_2}
    \chemical{GIVES} \chemical{2H_2O}
  \stopformula
\stopbuffer

\typebuffer

\getbuffer

The definition of the chemical part could be somewhat shorter:

\starttyping
\chemical{2H_2,PLUS,O_2,GIVES,2H_2O}
\stoptyping

or even:

\starttyping
\chemical{2H_2,+,O_2,->,2H_2O}
\stoptyping

A \TEX||addict will notice from these examples that the plus sign and the
arrow are on the baseline. Compare for example $+$ and \chemical{+}. In the
reaction you will see that the \chemical{+} and the \chemical{->} are
vertically aligned. 

You can use \type{PLUS}, \type{GIVES} and \type{EQUILIBRIUM} (\type{<->}) in
this command. With \type{MESOMERIC} or \type{<>} you will get
\chemical{MESOMERIC}. 

The reaction can be placed in the text. In that case a more compact display
is used: \chemical{2H_2,+,O_2,->,2H_2O}. Some finetuning with \type{\,} would
result in \chemical{2{\,}H_2,{\,}+{\,},O_2\,,->,\,2{\,}H_2O}. 

It is also possible to display bonds in textmode. For example if you want
\chemical{H,SINGLE,CH,DOUBLE,HC,SINGLE,H} you should type
\type{\chemical{H,SINGLE,CH,DOUBLE,HC,SINGLE,H}} or something like this
\type{\chemical{H,-,CH,--,HC,-,H}}. A triple bond can be defined as
\type{TRIPLE} or \type{---}: \chemical{HC,TRIPLE,CH}. 

We return two the display||mode. The second and third argument can be used to
add text to the reaction: 

\startbuffer
\placeformula
  \startformula
    \chemical{2H_2}{hydrogen} \chemical{PLUS} \chemical{O_2}{oxygen}
    \chemical{GIVES}{heat} \chemical{2H_2O}{water}
   \stopformula
\stopbuffer

\typebuffer

So we can also place text over and under symbols!

\getbuffer

The last argument is placed under the compound. 

\startbuffer
\placeformula
  \startformula
    \chemical{H_2O}{liquid}{water}
    \hbox{c.q.}
    \chemical{H_2O}{water}
  \stopformula
\stopbuffer

\getbuffer

The formula above is defined with:

\typebuffer

The size of the formulas or reactions in the running text can be set up 
with:

\placetable
{Set up in text formulas.}
\starttable[|lT|Tl|lT|]
\HL
\VL \bf parameter \VL \bf set up \VL \bf default \VL\SR
\HL
\VL size  \VL small medium big \VL big \VL\SR
\HL
\stoptable

The definition \type{\chemical{H,SINGLE,CH,DOUBLE,HC,SINGLE,H}} result with
\type{big}, \type{medium} and \type{small} in the following formulas: 

\leavevmode\hbox
  {\setupchemical[textsize=big]\chemical{H,-,CH,--,HC,-,H}\hskip2em
   \setupchemical[textsize=medium]\chemical{H,-,CH,--,HC,-,H}\hskip2em
   \setupchemical[textsize=small]\chemical{H,-,CH,--,HC,-,H}}

\chapter{Subscripts}

Sub- and superscripts are placed somewhat lower as is recommended by Knuth in
the {\TeX}Book. The rather strange chemical compound that is shown on
page~179 of the {\TeX}Book is defined with
\type{\chemical{Fe_2^{+2}Cr_2O_4}}. This will result in 
\chemical{Fe_2^{+2}Cr_2O_4}. Without correction it would have been: $\rm
Fe_2^{+2}Cr_2O_4$. 

The position of the subscript is determined by the parameter \type{offset}:
\type{HIGH} or \type{LOW}. This position can be influenced locally (per
substituent) as is shown in the example below. 

\startbuffer
\startchemical[frame=on,option=test,alternative=2]
  \chemical[SIX,B,R13,HIGH,RZ1,LOW,RZ3][NH_3,NH_3]
\stopchemical
\stopbuffer

\startfiguretext
  [left][]
  {none}
  {\getbuffer}
\voorbeeld[ex:test]
\typebuffer
\stopfiguretext

However, it is recommended to set up this parameter globally to obtain
consistent formulas. 

The values \type{LOW} and \type{HIGH} can also be used in text formulas. For
example if you type \type{\chemical{HIGH,H_2O}} then you will get 
\chemical{HIGH,H_2O} and \type{\chemical{LOW,H_2O}} will become
\chemical{LOW,H_2O}. 

\vervolgblad
  {Backgrounds}
  [left=1500,width=4000]
  [FIVE,ROT3,B125,+SB3,-SB4,ER35,SR4,CRZ35,Z4,RZ4,Z0]
  [O,O,N,CH_3]

\chapter{Installation}

The package \PPCHTEX\ is developed for use in combination with \CONTEXT.
\PPCHTEX\ is activated in \CONTEXT\ by:\footnote{The macros in file
\type{ppchtex.tex}, are loaded by typing \type{m-chemic.tex} (the~\type{m}
stands for module).} 

\starttyping
\usemodule[pictex,chemic]
\stoptyping

This command loads the \PICTEX\ macros so \PPCHTEX\ knows what output is
needed. The chemical macros are not automatically available. 

The package can be used in combination with \PLAIN\ \TEX\ or \LATEX. In that
case the file \type{m-chemie.sty} is also used. \PPCHTEX\ is then activated
by \type{\documentstyle}: 

\starttyping
\documentstyle[m-pictex,m-chemie]{}
\stoptyping

In \LATEX$2_{\textstyle\varepsilon}$ it is somewhat different:

\starttyping
\usepackage{m-pictex}
\usepackage{m-chemie}
\stoptyping

The file \type{m-pictex} takes care of an efficient loading of \PICTEX. This
is necessary because \LATEX\ allocates a lot of \type{\dimen}s. Because of
the userinterface a big version of \TEX\ is needed to run \PPCHTEX. 

\PPCHTEX\ can be used in three languages. The actual language can be
activated by: 

\startregelcorrectie
\starttable[|l|l|]
\HL 
\NC \bf language  \NC \bf files  \NC\SR 
\HL 
\NC dutch   \NC \type{m-che-nl} = \type{m-chemie.sty} \NC\FR
\NC english \NC \type{m-che-en} = \type{m-chemic.sty} \NC\MR
\NC german  \NC \type{m-che-de} = \type{m-chemie.sty} \NC\LR
\HL 
\stoptable
\stopregelcorrectie

In the file \type{ppchtex.noc} you can see the coupling between \CONTEXT\ and
macropackage. This file also loads the system modules of \CONTEXT. 

The total distribution consists of the definition files \type{ppchtex.tex} and
\type{ppchtex.noc} the starting files \type{m-che-nl.tex},
\type{m-che-en.tex} and \type{m-che-de.tex} and the alternative starting
files \type{m-chemie.tex} and \type{m-chemic.tex}. 

\chapter{Extensions}

Users of \PPCHTEX\ are free to use and alter the macros. However one should
be careful because most macros are stil under development. Some macros may
look more complex than necessary, but this is due to the userfriendly
interface of \CONTEXT\ and \PPCHTEX. 

Commands like \type{\setup...} make it possible to make readable
\ASCII||layouts. Compare for example: 

\start
\setuptyping[option=color]
\starttyping
\setupchemical[size=small]
\stoptyping
\stop

and:

\start
\setuptyping[option=color]
\starttyping
\setupchemical
  [size=small,
   scale=500,
   textsize=big]
\stoptyping
\stop

The set up can be defined in a random order and newlines and spaces are not
interpreted. 

Originally \PPCHTEX\ was meant to be a module in \CONTEXT, therefore the
package is multi lingual. 

If you study the file \type{ppchtex.tex} you may notice that
\type{\processaction} macros are being used while interpretating the keys in
\type{\chemical[ ]}. These kind of macros are relatively slow but then the
\PPCHTEX\ interface is very flexible. 

\chapter{Fonts}

The macros are in mathematical mode and therefore use \type{\textfont},
\type{\scriptfont} and \type{\scriptscriptfont}. When needed the
\type{\fontdimen}s~14, 16 and~17 of \type{\font2} are adapted. The size of
the actual font is derived from $x$ height (\type{\fontdimen5}). 

Changes in \type{\fontdimen} s have are global, so grouping makes no sense.
The dimensions are therefore continually set and reset. This solution may
seem poor but alternatives are not failsave and will result in problems with
scaled fonts. 

Typesetting atoms and molecules (text) during processing are rather time
consuming. Speed depends on the complexity of the macro \type{\rm}. 

\chapter{Definitions}

The interface of \PPCHTEX\ is derived from the \CONTEXT\ interface. This
means that the interface is multi langual. The advantage is that one can use
\PPCHTEX\ in his or her own language. The disadvantage is the fact that
macros have to be shared between languages. 

At this moment the \CONTEXT\ commands and parameters are dutch. 
\PPCHTEX\ however has english commands. 

\starttyping
\startchemical
  \chemical[SIX,B,C]
\stopchemical
\stoptyping

Set ups are more difficult. We use system constants and variables that are 
dutch. In due time these will be translated. Adaption of these 
constants and variables is not too difficult. 

\starttyping
\setupchemical[\c!breed=10cm,\c!height=\v!passend]
\stoptyping

Parameters are preceded by \type{\c!} and set ups by \type{\v!}. This only
works when \type{!} is a character. That is the reason that these set ups
have to be surrounded by \type{\unprotect} and \type{\protect}. For example: 

\starttyping
\unprotect

\setupchemical[\c!width=10cm,\c!height=\v!passend]

\startchemical
  \chemical[SIX,B,C]
\stopchemical

\protect
\stoptyping

More information on the interface can be found in the documentation of the
\CONTEXT\ modules from the \type{mult} group. 

\chapter{Color}

In \CONTEXT\ you can colorize parts of a structure. 
In~\in{example}[ex:color] the substituent as well as the bond
are colored red. 

\startbuffer
\startchemical[axis=on,frame=on,width=5000]
  \chemical[SIX,B]
  \color[red]{\chemical[SIX,R1,RZ1][COOH]}
\stopchemical
\stopbuffer

\startfiguretext
  [left][]
  {none}
  {\getbuffer}
\voorbeeld[ex:color]
\typebuffer
\stopfiguretext

First the color mechanism has to be activated by
\type{\setupcolors[state=start]}. 

\chapter[txt:cooh]{Interaction}

In combination with \CONTEXT\ \PPCHTEX\ supports interactive texts. An
interactive text is a text that can be consulted on a computerscreen and
contains many hyperlinked textareas. This means that clicking on such an area
will result in a jump to the target area. 

\startbuffer
\startchemical[axis=on,frame=on,width=5000]
  \chemical[SIX,B]
  \chemical[sub:cooh][SIX,R1,RZ1][COOH]
\stopchemical
\stopbuffer

\startfiguretext
  [left][]
  {none}
  {\getbuffer}
\voorbeeld[ex:interaction]
\typebuffer
\stopfiguretext

We see a new argument: the reference \type{[sub:cooh]}. This means that we
can refer from the text \goto{\chemical{COOH}}[sub:cooh] to the structure 
with: 

\starttyping
... text ... \goto{\chemical{COOH}}[sub:cooh] ... text ...
\stoptyping

In this definition \type{\goto} is a \CONTEXT||command. 
We can also refer from the structure to a particular part of the text. 

Clicking in \chemical{COOH} in the structure results in a jump to the text
that is marked with: 

\starttyping
\paragraph[txt:cooh]{Substituents}

... text ... \chemical{COOH} ... text ... 
\stoptyping 

A combination is also possible. In that case it is necessary to mark the
reference with \type{\chemical} and to refer in the text with
\type{\gotochemical}. 

The coupling of the interaction mechanism is done with macros:

\starttyping
\localgotochemical   {reference} {text}
\localthisischemical {reference}
\stoptyping

You can see the \CONTEXT\ sources for more information. 

\vervolgblad
  {Overview}
  [left=1250,width=4000,bottom=1500,height=4000]
  [SIX,B,C,MOV1,B,ER6,CRZ6]
  [O]

\input man-comm.tex

\stoptext
